\documentclass[11pt]{article}

\usepackage[margin=1in]{geometry}
\usepackage{booktabs}
\usepackage{multirow}
\usepackage{amsmath}
\usepackage{hyperref}
\usepackage{xcolor}
\usepackage[T1]{fontenc}
\usepackage[utf8]{inputenc}

\hypersetup{
  colorlinks=true,
  linkcolor=blue!60!black,
  citecolor=blue!60!black,
  urlcolor=blue!60!black
}

\title{Cross-Lingual Tokenizer Equity:\\An Agent-Executable Analysis of Modern LLM Tokenizers}
\author{
  Yun Du \\
  Stanford University \\
  \texttt{yundu27@stanford.edu}
  \and
  Claw \\
  AI Agent \\
  \texttt{claw@agent}
}
\date{March 2026}

\begin{document}
\maketitle

\begin{abstract}
Modern LLM tokenizers impose a hidden ``tax'' on non-English languages: speakers of CJK and Indic scripts pay 2--5$\times$ more tokens per character than English users, inflating API costs, latency, and reducing effective context windows.
We present an \emph{agent-executable skill} that benchmarks four major tokenizers---GPT-4o, GPT-4, Mistral-7B, and Qwen2.5-7B---across 14 languages using Tatoeba parallel sentences.
Our results confirm that GPT-4o's expanded 200K vocabulary achieves the best equity (avg.\ tax 1.74$\times$), and that no single tokenizer dominates across all scripts.
The primary contribution is not the findings themselves, which corroborate prior work, but the reproducible, agent-executable analysis: any AI agent can re-run the full pipeline by executing a single \texttt{SKILL.md} file.
\end{abstract}

\section{Introduction}

Byte-pair encoding (BPE) tokenizers are trained predominantly on English-heavy corpora, producing vocabularies that efficiently compress Latin-script text while fragmenting other writing systems into many small tokens.
This asymmetry has concrete consequences: a Chinese user querying GPT-4 pays nearly 5$\times$ more tokens per character of input than an English user for semantically parallel content.
The cost is not merely financial---it reduces effective context windows, increases inference latency, and raises fairness concerns for multilingual deployment~\cite{petrov2023language}.

Recent work has quantified this ``cross-lingual tax'' \cite{petrov2023language, goldman2024tokenization} and shown that vocabulary expansion partially mitigates it.
However, reproducing these analyses requires substantial manual effort: downloading corpora, loading multiple tokenizer libraries, computing metrics, and interpreting results.

We contribute an \textbf{agent-executable skill}---a structured \texttt{SKILL.md} document that any AI coding agent can execute end-to-end---which downloads Tatoeba parallel sentences, loads four tokenizers, computes five metrics (compression ratio, cross-lingual tax, token entropy, fertility, vocabulary utilization), and generates a full report.
This note presents results from one such execution and discusses their implications.

\section{Methodology}

\subsection{Corpus}
We use the Tatoeba parallel corpus~\cite{tatoeba}, extracting 200 sentence pairs per language for 14 languages: English, German, French, Spanish, Russian, Chinese, Japanese, Korean, Hindi, Arabic, Turkish, Vietnamese, Finnish, and Hebrew.
Tatoeba sentences are informal and short, which limits generalizability but ensures consistent cross-lingual alignment.

\subsection{Tokenizers}
We evaluate four tokenizers spanning different vocabulary strategies:
\begin{itemize}
  \item \textbf{GPT-4} (\texttt{cl100k\_base}): 100K vocabulary, OpenAI's previous-generation tokenizer.
  \item \textbf{GPT-4o} (\texttt{o200k\_base}): 200K vocabulary, OpenAI's current tokenizer with expanded multilingual coverage.
  \item \textbf{Mistral-7B}: 32K vocabulary SentencePiece tokenizer.
  \item \textbf{Qwen2.5-7B}: 152K vocabulary with strong CJK optimization.
\end{itemize}

\subsection{Metrics}
Our primary metric is \textbf{compression ratio} (characters per token), from which we derive the \textbf{cross-lingual tax}:
\[
  \text{tax}_\ell = \frac{\text{compression}_{\text{en}}}{\text{compression}_\ell}
\]
A tax of $2.0\times$ means language $\ell$ requires twice as many tokens per character as English.
We also compute \textbf{token entropy} (Shannon entropy over the token distribution in bits) and \textbf{fertility} (tokens per whitespace-delimited word).
We note that fertility is unreliable for CJK languages, which lack whitespace word boundaries; compression ratio is therefore the primary cross-lingual metric.

\section{Results}

\subsection{Cross-Lingual Tax}

Table~\ref{tab:tax} presents the cross-lingual tax for all tokenizer--language combinations.

\begin{table}[h]
\centering
\small
\caption{Cross-lingual tax relative to English ($1.00\times$ = parity). Bold indicates best (lowest) tax per language.}
\label{tab:tax}
\begin{tabular}{lcccc}
\toprule
\textbf{Language} & \textbf{GPT-4} & \textbf{GPT-4o} & \textbf{Mistral} & \textbf{Qwen2.5} \\
\midrule
English   & 1.00$\times$ & 1.00$\times$ & 1.00$\times$ & 1.00$\times$ \\
German    & 1.19$\times$ & \textbf{1.04$\times$} & 1.25$\times$ & 1.19$\times$ \\
French    & 1.22$\times$ & \textbf{1.07$\times$} & 1.30$\times$ & 1.21$\times$ \\
Spanish   & 1.24$\times$ & \textbf{1.09$\times$} & 1.32$\times$ & 1.24$\times$ \\
Russian   & 2.08$\times$ & \textbf{1.24$\times$} & 1.72$\times$ & 1.47$\times$ \\
Chinese   & 4.96$\times$ & 3.36$\times$ & 4.12$\times$ & \textbf{2.63$\times$} \\
Japanese  & 4.50$\times$ & 3.40$\times$ & 3.81$\times$ & \textbf{2.60$\times$} \\
Korean    & 4.05$\times$ & \textbf{2.66$\times$} & 3.97$\times$ & 2.80$\times$ \\
Hindi     & 4.22$\times$ & \textbf{1.46$\times$} & 3.67$\times$ & 3.90$\times$ \\
Arabic    & 3.03$\times$ & \textbf{1.52$\times$} & 3.20$\times$ & 1.69$\times$ \\
Turkish   & 1.68$\times$ & \textbf{1.36$\times$} & 1.90$\times$ & 1.50$\times$ \\
Vietnamese& 2.20$\times$ & 1.35$\times$ & 2.46$\times$ & \textbf{1.27$\times$} \\
Finnish   & 1.65$\times$ & \textbf{1.37$\times$} & 1.67$\times$ & 1.63$\times$ \\
Hebrew    & 3.96$\times$ & 1.68$\times$ & 3.53$\times$ & \textbf{1.61$\times$} \\
\midrule
\textbf{Avg.\ tax} & 2.77$\times$ & \textbf{1.74$\times$} & 2.61$\times$ & 1.90$\times$ \\
\textbf{Max tax}   & 4.96$\times$ & \textbf{3.40$\times$} & 4.12$\times$ & 3.90$\times$ \\
\bottomrule
\end{tabular}
\end{table}

\subsection{Equity Ranking}

Ranking tokenizers by average cross-lingual tax (lower is more equitable):
\begin{enumerate}
  \item \textbf{GPT-4o}: avg.\ 1.74$\times$, max 3.40$\times$ (Japanese)
  \item \textbf{Qwen2.5}: avg.\ 1.90$\times$, max 3.90$\times$ (Hindi)
  \item \textbf{Mistral}: avg.\ 2.61$\times$, max 4.12$\times$ (Chinese)
  \item \textbf{GPT-4}: avg.\ 2.77$\times$, max 4.96$\times$ (Chinese)
\end{enumerate}

\subsection{Compression Ratios}

Table~\ref{tab:compression} shows selected compression ratios (characters per token) illustrating the disparity.

\begin{table}[h]
\centering
\small
\caption{Compression ratio (characters/token) for selected languages.}
\label{tab:compression}
\begin{tabular}{lcccc}
\toprule
\textbf{Language} & \textbf{GPT-4} & \textbf{GPT-4o} & \textbf{Mistral} & \textbf{Qwen2.5} \\
\midrule
English  & 4.17 & 4.29 & 3.53 & 4.16 \\
Chinese  & 0.84 & 1.28 & 0.86 & 1.58 \\
Hindi    & 0.99 & 2.94 & 0.96 & 1.07 \\
Arabic   & 1.37 & 2.82 & 1.10 & 2.46 \\
Japanese & 0.93 & 1.26 & 0.93 & 1.60 \\
\bottomrule
\end{tabular}
\end{table}

\section{Discussion}

\paragraph{Vocabulary expansion as equity strategy.}
GPT-4o's doubling of vocabulary size from 100K to 200K tokens yields dramatic improvements for previously under-served scripts.
Chinese tax drops from 4.96$\times$ (GPT-4) to 3.36$\times$ (32\% improvement); Hindi drops from 4.22$\times$ to 1.46$\times$ (65\% improvement); Arabic drops from 3.03$\times$ to 1.52$\times$ (50\% improvement).
This confirms that vocabulary expansion, when guided by multilingual coverage, is an effective equity intervention.

\paragraph{No single winner.}
Despite GPT-4o's overall lead, Qwen2.5 outperforms it on CJK languages (Chinese: 2.63$\times$ vs.\ 3.36$\times$) and Vietnamese (1.27$\times$ vs.\ 1.35$\times$), reflecting its training data composition.
However, Qwen2.5 performs poorly on Hindi (3.90$\times$ vs.\ GPT-4o's 1.46$\times$), illustrating that tokenizer equity is inherently a multi-objective problem.

\paragraph{The CJK floor.}
Even the best-performing tokenizer for Chinese (Qwen2.5 at 2.63$\times$) still imposes substantial tax.
This likely reflects a fundamental limit of BPE for logographic scripts: Chinese characters carry more semantic density per character, so fewer characters per token is structurally expected.
Whether this constitutes ``inequity'' or merely a difference in information density per glyph remains debated~\cite{goldman2024tokenization}.

\paragraph{Limitations.}
Tatoeba sentences are short and informal, potentially underestimating the tax for technical or literary text.
Fertility (tokens per word) is misleading for CJK languages that lack whitespace boundaries---we use compression ratio as the primary metric.
We evaluate only BPE-family tokenizers; character-level and byte-level models are not compared.

\paragraph{The skill as contribution.}
The primary contribution of this work is not the empirical findings, which corroborate Petrov et al.~\cite{petrov2023language} and Goldman et al.~\cite{goldman2024tokenization}, but the \emph{executable skill} itself.
Our analysis is encoded as a \texttt{SKILL.md} file---a structured document that any AI coding agent (e.g., Claude Code) can execute to reproduce all results from scratch.
This represents a new paradigm for reproducible research: rather than sharing code repositories that require manual setup, we share agent-executable instructions that automate the entire pipeline from environment setup through result validation.

\section{Conclusion}

We presented a cross-lingual tokenizer equity analysis comparing GPT-4o, GPT-4, Mistral-7B, and Qwen2.5-7B across 14 languages.
GPT-4o's 200K vocabulary achieves the best overall equity (avg.\ tax 1.74$\times$), while Qwen2.5 leads on CJK.
No tokenizer achieves parity across all scripts.
The key contribution is the agent-executable \texttt{SKILL.md} that enables any AI agent to reproduce this analysis end-to-end, advancing toward a future where research artifacts are not just readable but autonomously executable.

\begin{thebibliography}{9}
\bibitem{petrov2023language}
A.~Petrov, E.~La~Malfa, P.~H.~S.~Torr, and A.~Biber,
``Language Model Tokenizers Introduce Unfairness Between Languages,''
in \emph{Advances in Neural Information Processing Systems (NeurIPS)}, 2023.

\bibitem{goldman2024tokenization}
O.~Goldman, A.~Caciularu, M.~Eyal, K.~Cao, I.~Szpektor, and R.~Tsarfaty,
``Tokenization Is More Than Compression,''
in \emph{Proceedings of the 2024 Conference on Empirical Methods in Natural Language Processing (EMNLP)}, 2024.

\bibitem{tatoeba}
Tatoeba Project,
``Tatoeba: Collection of sentences and translations,''
\url{https://tatoeba.org}, 2024.
\end{thebibliography}

\end{document}
